% GENERATED FILE. Edit canonical lessons, then run npm run generate:books.
% canonical-source-hash: fnv1a64:002a7edb9e61f4bd
% canonical-lessons: ML-C112-bank, ML-C112-poleesu, ML-C112-aashupathri, ML-R112-places-recall

\chapter{Where to Go}
\label{ch:ml-places}
\hlchaptermodality{112}

\begin{chapteropening}
\textbf{By the end of this chapter:} \emph{I can name a bank, a hospital and the police, and tell which everyday words my English ear can decode for me.}

From cold: three places, and which of the three the learner's English ear could not reach.
\end{chapteropening}

\section[bāṅkŭ]{\ml{ബാങ്ക്} (bāṅkŭ) — a bank}

\label{lesson:ML-C112-bank}



\begin{quote}
Say \emph{money}. Say \emph{a telephone}. One of those two you learned as an English word wearing Malayalam letters. Hold on to that feeling — this chapter hands you two more of them.
\end{quote}

\subsection*{What to know first}
\textbf{\ml{ബാങ്ക്}} (\emph{bāṅkŭ}) — \textbf{a bank}.

You have been able to say \textbf{\ml{പണം}} for a while. Now you can say where it lives.

\begin{culture}
Read it aloud and you can hear that nothing was done to it. English \emph{bank} came into Malayalam whole, and the letters spell out the English sounds one for one.

That is the second time you have seen this. The word for a telephone arrived the same way, and for the same reason: the thing itself arrived from outside, and the name came along in the box with it. A language that meets a new object often takes the object's name rather than building a fresh one.

Which means you can read \textbf{\ml{ബാങ്ക്}} before anybody teaches it to you — and that is a real skill, not a trick. There are a great many words like these two.
\end{culture}

\subsection*{Your turn}
Say these aloud:

\begin{itemize}
\raggedright
  \item \emph{bāṅkŭ}
  \item the money, then the bank
  \item the other English word you already read this way
\end{itemize}

\begin{tcolorbox}[breakable,colback=teal!4,colframe=teal!35!black,title={Before you move on}]
A bank? (\textbf{\ml{ബാങ്ക്}}.) What is kept there? (\textbf{\ml{പണം}}.) Did Malayalam build this word or borrow it? (\textbf{Borrowed} — you can hear the English in it.)
\end{tcolorbox}
\section[pōlīsŭ]{\ml{പോലീസ്} (pōlīsŭ) — the police}

\label{lesson:ML-C112-poleesu}



\begin{quote}
Say \emph{a bank}. Say \emph{to help}. You have had the verb for helping since long before this chapter; you have had the bank for about a minute.
\end{quote}

\subsection*{What to know first}
\textbf{\ml{പോലീസ്}} (\emph{pōlīsŭ}) — \textbf{the police}.

One word covers the force and its officers alike, so you do not need a second word to point at a person.

\begin{culture}
This is the same story as the bank, one word later. English \emph{police} came in whole, and the letters trace the English sounds.

Two borrowings in two lessons is not a coincidence. Both name something that came to Kerala as an institution with paperwork and a building, and a thing that arrives already named tends to keep the name it came with.

So the pattern now has two members you can point at, and you should expect more. When a word looks long and sounds oddly English as you sound it out, trust that ear before you reach for a dictionary.
\end{culture}

\subsection*{Your turn}
Say these aloud:

\begin{itemize}
\raggedright
  \item \emph{pōlīsŭ}
  \item the bank, then the police
  \item the police, then the verb for helping
\end{itemize}

\begin{tcolorbox}[breakable,colback=teal!4,colframe=teal!35!black,title={Before you move on}]
The police? (\textbf{\ml{പോലീസ്}}.) A bank? (\textbf{\ml{ബാങ്ക്}}.) What do these two words have in common? (\textbf{Both were borrowed whole from English.})
\end{tcolorbox}
\section[āśupatri]{\ml{ആശുപത്രി} (āśupatri) — a hospital}

\label{lesson:ML-C112-aashupathri}



\begin{quote}
Say \emph{a doctor}. Say \emph{medicine}. You have had the person and the substance for some time. You have never had the building they are both in.
\end{quote}

\subsection*{What to know first}
\textbf{\ml{ആശുപത്രി}} (\emph{āśupatri}) — \textbf{a hospital}.

The doctor, the medicine, the building. The set is complete now.

\begin{culture}
Sound this one out and notice what does \textbf{not} happen.

The last two words handed themselves to you. You met \textbf{\ml{ബാങ്ക്}}, heard English \emph{bank}, and were done. \textbf{\ml{പോലീസ്}} did the same. Your English ear did the work.

\textbf{\ml{ആശുപത്രി}} will not do that for you. Whatever its history, the word that reaches you today has Malayalam's own shape and gives your English ear nothing to catch. This lesson is not going to tell you where it came from, because that question has more than one answer in circulation and a guess dressed up as a fact would be worth less to you than the honest gap.

What matters for reading is the part that is certain: some words you can decode from English and some you cannot, and you will not know which until you try. So try first, and learn the word properly when the ear comes back empty.
\end{culture}

\subsection*{Your turn}
Say these aloud:

\begin{itemize}
\raggedright
  \item \emph{āśupatri}
  \item the doctor, the medicine, the hospital
  \item which of this chapter's three words your English ear could not catch
\end{itemize}

\begin{tcolorbox}[breakable,colback=teal!4,colframe=teal!35!black,title={Before you move on}]
A hospital? (\textbf{\ml{ആശുപത്രി}}.) Who works there? (\textbf{\ml{വൈദ്യൻ}}.) Could you have read this word off an English one? (\textbf{No} — this one had to be learned.)
\end{tcolorbox}
\section[bāṅkŭ, pōlīsŭ, āśupatri]{(bank, police, hospital) — from cold}

\label{lesson:ML-R112-places-recall}



\begin{quote}
Close the lessons before this one. Say \emph{a bank}, \emph{the police} and \emph{a hospital}.
\end{quote}

\begin{grammarlens}[title={Grammar Lens: two your ear could read, one it could not}]
\noindent
\begin{tabularx}{\linewidth}{@{}>{\raggedright\arraybackslash}X>{\raggedright\arraybackslash}X>{\raggedright\arraybackslash}X@{}}
\toprule
\textbf{} & \textbf{} & \textbf{} \\
\midrule
\textbf{\ml{ബാങ്ക്}} & \emph{bāṅkŭ} & English, audibly \\
\textbf{\ml{പോലീസ്}} & \emph{pōlīsŭ} & English, audibly \\
\textbf{\ml{ആശുപത്രി}} & \emph{āśupatri} & gives the English ear nothing \\
\bottomrule
\end{tabularx}

\textbf{Two of the three you could have decoded before anybody taught them.} Do not read a rule into it — the count here is two against one, and the next three words you meet could fall the other way. What this chapter gives you is the habit, not a law: sound a long word out before you look it up, because sometimes the word answers back in a language you already have.

And when it does not answer, that is information too. \textbf{\ml{ആശുപത്രി}} told you quickly that it was going to have to be learned.
\end{grammarlens}

\begin{grammarlens}[title={What you've built}]
\textbf{You could already name what you needed; you could not name where to go for it.} \textbf{\ml{പണം}} was yours from the shopping words, \textbf{\ml{വൈദ്യൻ}} from the words for what people do, \textbf{\ml{മരുന്ന്}} from the chapter on falling ill, and the verb for helping from long before any of them.

Every one of those was a thing or a person. This chapter adds the addresses. Money goes to a bank, illness goes to a hospital, and when you need help of the kind a neighbour cannot give, there is the police.

That is the difference between knowing words and being able to get through a day.
\end{grammarlens}

\subsection*{Your turn}
Say these aloud:

\begin{itemize}
\raggedright
  \item \emph{bāṅkŭ}, \emph{pōlīsŭ}, \emph{āśupatri}
  \item which of the three your English ear could not catch
  \item the money and the place it goes
  \item the doctor, the medicine, and the building that holds both
\end{itemize}

\begin{tcolorbox}[breakable,colback=teal!4,colframe=teal!35!black,title={Before you move on}]
A bank, the police, a hospital? (\textbf{\ml{ബാങ്ക്}, \ml{പോലീസ്}, \ml{ആശുപത്രി}}.) Which one had to be learned rather than decoded? (\textbf{\ml{ആശുപത്രി}}.) What should you do with a long word before reaching for a dictionary? (\textbf{Sound it out} — it may answer in English.)
\end{tcolorbox}
